\documentclass[UTF8,a4paper,11pt]{ctexart}
\usepackage[margin=22mm]{geometry}
\usepackage{booktabs,amsmath,graphicx}
\title{前期工作：指标表与实验依据}
\author{}
\date{2026年10月26日}
\begin{document}
\maketitle
\section*{已完成的基线与方法验证}
% Source: audited GOAL/OBJECT seeds 41,42,43; units percent; sample SD.
\begin{table}[htbp]
\centering
\caption{MLR+IFA复现与回放方法的已完成闭环结果}
\label{tab:preliminary-baselines}
\begin{tabular}{llcc}
\toprule
套件 & 方法 & AUC（\%） & 最终成功率（\%） \\
\midrule
GOAL & 原MLR+IFA & 76.35 $\pm$ 6.38 & 76.67 $\pm$ 9.21 \\
GOAL & MLR+IFA+RA & 81.98 $\pm$ 3.00 & 80.83 $\pm$ 5.05 \\
OBJECT & 原MLR+IFA & 68.30 $\pm$ 3.12 & 65.42 $\pm$ 5.20 \\
OBJECT & MLR+IFA+RA & 73.23 $\pm$ 1.78 & 69.58 $\pm$ 5.64 \\
OBJECT & MLR+IFA+RA+ALIGN & 80.17 $\pm$ 3.66 & 75.83 $\pm$ 3.82 \\
\bottomrule
\end{tabular}
\par\smallskip
{\footnotesize seed41、42、43的均值$\pm$样本标准差，每个阶段--任务单元20回合。
最终成功率为最后阶段四任务平均，数值为项目复现结果。}
\end{table}

% S41 accepted final evaluation, separate from three-seed means.
\begin{table}[htbp]
\centering
\caption{S方法在GOAL seed41上的最终闭环结果}
\label{tab:preliminary-s41}
\begin{tabular}{lrrrrrr}
\toprule
方法 & Task6 & Task7 & Task8 & Task9 & 旧任务平均 & 全部平均 \\
\midrule
S & 45.00 & 100.00 & 100.00 & 95.00 & 81.67 & 85.00 \\
\bottomrule
\end{tabular}
\par\smallskip
{\footnotesize 单位为\%，每任务20正式回合。另有40次训练采集，单种子初步验证。}
\end{table}

\paragraph{指标定义}
令$r_{s,t}$表示学完阶段$s$后在任务$t$上的闭环成功率，$T=4$，任务索引重编号为$1,\ldots,T$。
\begin{align}
\mathrm{AUC}&=\frac{1}{T}\sum_{t=1}^{T}\frac{1}{T-t+1}\sum_{s=t}^{T}r_{s,t},\\
\mathrm{OldFinal}&=\frac{1}{T-1}\sum_{t=1}^{T-1}r_{T,t},\\
\mathrm{FinalAll}&=\frac{1}{T}\sum_{t=1}^{T}r_{T,t}.
\end{align}
表中数值为上述比例乘以100。三种子的离散程度使用样本标准差（分母$n-1$）。

\clearpage
\section*{真实训练记录}
\begin{figure}[htbp]
\centering\includegraphics[width=\linewidth]{images/training_curves_S41.pdf}
\caption{S方法在GOAL seed41上的实际训练曲线。记录总量除以每轮batch数，未平滑。}
\end{figure}
S完成Task6至Task9各100轮训练，共400轮、76\,100次主更新。epoch从0编号，第99轮对应阶段的第100个完整训练轮次。
\begin{verbatim}
arm: S      stage: 9       epoch: 99
steps: 76100              batches: 337
valid: 3218               invalid: 147
lr_after: 0.0
\end{verbatim}
\section*{真实双相机关键帧}
\begin{figure}[htbp]
\centering
\begin{tabular}{ccc}
第0步 & 第65步 & 第77步\\
\includegraphics[width=.26\linewidth]{images/first_mean_delta_agentview_image.png}&\includegraphics[width=.26\linewidth]{images/first_mode_switch_agentview_image.png}&\includegraphics[width=.26\linewidth]{images/terminal_agentview_image.png}\\
\includegraphics[width=.26\linewidth]{images/first_mean_delta_robot0_eye_in_hand_image.png}&\includegraphics[width=.26\linewidth]{images/first_mode_switch_robot0_eye_in_hand_image.png}&\includegraphics[width=.26\linewidth]{images/terminal_robot0_eye_in_hand_image.png}
\end{tabular}
\caption{已有ALIGN诊断rollout的成功示例：GOAL Task8（将碗放到盘子上），seed41，Stage9，初态3。上行为外部相机，下行为腕部相机。图片是已保存的稀疏关键帧。}
\end{figure}
\end{document}
